% ---------------------------------------------------------------------
%  CHƯƠNG 6 — KẾT LUẬN VÀ HƯỚNG PHÁT TRIỂN
% ---------------------------------------------------------------------
\chapter{Kết luận và hướng phát triển}
\label{ch:ketluan}

\begin{muctieu}
Chương này tổng kết những gì đồ án đã làm được, rút ra bài học và đề xuất
hướng phát triển tiếp theo.
\end{muctieu}

\section{Kết luận}

Đồ án đã xây dựng một hệ thống trí tuệ nhân tạo gợi ý đường đi của cơn bão
ảnh hưởng Việt Nam, đi trọn từ dữ liệu IBTrACS 1980--2024, tạo đặc trưng,
huấn luyện và so sánh bảy phương pháp, đến triển khai API và bản đồ web. Bốn
trên năm mục tiêu đạt đầy đủ (Bảng~\ref{tab:muctieu}): các phương pháp học giảm
sai số 72 giờ khoảng 200~km so với persistence với ý nghĩa thống kê, và nón bất
định phủ 82--86\% vị trí thực. Mục tiêu triển khai mới đạt một phần vì mô hình
tốt nhất, có dòng dẫn đường, chưa được đưa vào phục vụ.

Về mặt huấn luyện, đồ án rút ra ba kết quả chính: dự báo độ dịch chuyển tốt hơn
hẳn dự báo tọa độ tuyệt đối ở hạn ngắn; với dữ liệu chỉ có quỹ đạo, LSTM,
Transformer, k-NN và gradient boosting tương đương nhau; và bổ sung thông tin
dòng dẫn đường khí quyển là thay đổi duy nhất cải thiện có ý nghĩa (giảm 35~km
tại 72 giờ).

\section{Bài học rút ra}

Bài học lớn nhất là \textbf{thông tin đầu vào quan trọng hơn kiến trúc mô hình}.
Chúng em đã thử nhiều cách làm mạng phức tạp hơn mà không cải thiện; chỉ khi
đưa vào dữ liệu phản ánh đúng nguyên nhân vật lý của chuyển động bão thì sai số
mới giảm. Nếu làm lại, chúng em sẽ tìm hiểu về khí tượng trước để chọn đặc trưng, thay vì
bắt đầu từ việc thử kiến trúc.

Bài học thứ hai là về \textbf{kỷ luật đánh giá}: chia tập theo năm, chọn mô hình
chỉ trên validation, chạy nhiều seed và kiểm định bootstrap. Nhiều khác biệt vài
chục km trông có vẻ đáng kể hoá ra không có ý nghĩa thống kê. Chúng em cũng phải
bỏ qua một cấu hình có điểm test tốt nhất vì nó không được chọn trên validation.

Về vai trò của trí tuệ nhân tạo: AI học được các kiểu quỹ đạo phổ biến và giảm
sai số rõ rệt so với ngoại suy đơn giản, nhưng vẫn thất bại ở các cơn bão đi
khác thường như YAGI, đúng những trường hợp nguy hiểm nhất. Vì vậy hệ thống chỉ
nên là công cụ gợi ý bổ sung, có giải thích (bão tương tự) và có vùng bất định,
không thay thế bản tin của cơ quan khí tượng có thẩm quyền.

\section{Hướng phát triển}

\begin{itemize}[leftmargin=1.6em, itemsep=2pt]
  \item Đưa mô hình kết hợp có dòng dẫn đường vào vận hành, lấy đặc trưng từ
        phân tích GFS tại thời điểm dự báo, khắc phục hạn chế mô hình đang phục
        vụ chỉ dùng quỹ đạo.
  \item Dùng tái phân tích ERA5 độ phân giải 0,25° và đặc trưng dạng ảnh (mạng
        CNN trên trường gió) để nhận diện áp cao cận nhiệt, nhắm vào các ca khó
        như YAGI.
  \item Thay nón tròn bằng nón bất đối xứng theo hướng dọc và ngang quỹ đạo,
        hiệu chỉnh riêng theo cường độ bão.
  \item Kiểm định chéo theo nhiều cách chia năm và theo dõi sai số trên các cơn
        bão mới 2025--2026.
\end{itemize}
